\section{Introduction}\label{sec:intro}

Kubernetes has become the de-facto platform for orchestrating containerised applications. Organisations increasingly rely on it to deploy, scale and operate production software \citep{rahman2023k8smisconfig,shamim2020xicommandments}. The behaviour of a Kubernetes workload is not defined in application code but in declarative configuration files, known as \emph{manifests}. Among other things, they specify how many replicas run and how container health is probed. They also specify how much CPU and memory each container may use, which privileges it holds and how it is exposed on the network. These settings determine both the security posture and the availability of the running system. Configuration mistakes can give containers root or host-level privileges, leak secrets, allow resource exhaustion, or leave a service unavailable when a pod fails \citep{rahman2023k8smisconfig,zhang2026k8sconfigdefects,nsacisa2022k8shardening}.

Empirical studies show that human-written manifests often contain such problems. \citet{rahman2023k8smisconfig} found 1,051 security misconfigurations in 2,039 manifests from 92 open-source repositories. They grouped them into 11 categories, such as absent resource limits, absent security contexts and privileged containers. \citet{zhang2026k8sconfigdefects} later listed 15 categories of Kubernetes configuration defects, including probing, resource and scheduling defects. They found that eight widely used static analysis tools detect only 8 of these 15 categories. Meanwhile, large language models (LLMs) have become everyday assistants for software and infrastructure engineering. Examples are ChatGPT, Claude and Gemini, and open-weight models such as Qwen and DeepSeek. Developers already ask them for Kubernetes manifests: \citet{zhang2024genaismells} found 98 such requests in a corpus of developer conversations with ChatGPT. Research on general-purpose code has repeatedly shown that LLM assistants reproduce insecure patterns. About 40\% of programs generated by GitHub Copilot in security-relevant scenarios were vulnerable \citep{pearce2022asleep}. About half of the functionally correct backends generated by state-of-the-art models were exploitable \citep{vero2025baxbench}. For infrastructure-as-code, studies have found that syntactic validity and deployability are largely independent of security compliance \citep{santosvargas2026secureterraform,zhang2026iacgen}.

However, evidence on the quality of LLM-generated Kubernetes configurations is thin. The first study on the topic analysed 98 ChatGPT-generated manifests collected from developer conversations. It found that 35.8\% contained unset CPU and memory requirements \citep{zhang2024genaismells}. That study used uncontrolled prompts, a single model from 2023 and a single static analysis tool, and no human baseline. A later comparison with human-written manifests \citep{jcomsec2026aivshumank8s} re-used those 2023 measurements against an unrelated human corpus. As its authors state, it did not generate new manifests, used a single tool and executed nothing. Benchmarks for LLM-generated cloud configuration, such as CloudEval-YAML \citep{xu2024cloudevalyaml} and FDE-Bench \citep{ding2026fdebench}, do execute their outputs. But they use curated tasks with reference solutions rather than real projects, and do not compare against the manifests developers actually ship. Other Kubernetes-specific LLM research uses models to \emph{detect}, \emph{repair} or \emph{harden} existing manifests, not to write them \citep{malul2024genkubesec,ye2025llmsecconfig,sgancohen2025kubeguard}. Three questions therefore remain open. (i) How do configurations produced by current LLMs compare with the human-written configurations of \emph{the same applications}? (ii) How do different LLMs compare with each other under identical inputs? (iii) Do these configurations actually deploy and work on a real cluster, and does static analysis predict which ones do? They correspond to RQ1, RQ3 and RQ4 below. RQ2 describes the misconfigurations themselves, and RQ5 tests whether simply asking for production-ready output changes the answers.

This study evaluates LLM-generated Kubernetes configurations against the manifests that the same open-source projects ship. It addresses five research questions:
\begin{itemize}
    \item \textbf{RQ1:} How do LLM-generated Kubernetes configurations compare with the human-written configurations shipped by the same projects in correctness, security, reliability, requirement satisfaction and runtime behaviour?
    \item \textbf{RQ2:} Which misconfigurations occur most frequently in LLM-generated configurations?
    \item \textbf{RQ3:} How do different LLMs differ in the problems they introduce, and how consistent is each model across repeated generations?
    \item \textbf{RQ4:} Why do generated configurations fail at runtime, and do static analysis findings predict which ones work?
    \item \textbf{RQ5:} Does explicitly requesting production-ready output change the security, reliability and deployability of LLM-generated configurations?
\end{itemize}

Seven non-trivial open-source applications whose maintainers ship Kubernetes manifests are used. They range from stateful database-backed applications and multi-service systems to a cluster add-on. Each model receives the same \emph{docs bundle} per application. It contains the project's own README, Dockerfile and Docker Compose file at a pinned commit, with Kubernetes-specific passages redacted. So no researcher- or LLM-written specification sits between the project and the model. Five LLMs (GPT, Claude, Gemini, Qwen and DeepSeek) generate three samples per project under a neutral prompt and under a production-oriented prompt. Neutral-prompt outputs are compared with the shipped manifests. The production-oriented prompt measures the effect of the request itself. All 198 generated and 7 human-written configurations go through the same six-dimension benchmark: correctness (D1), security (D2), static reliability (D3), requirement satisfaction (D4), deployment with a functional smoke test on a disposable cluster (D5) and recovery from injected pod failures (D6).

The main contributions of this paper are as follows:
\begin{itemize}
    \item A \emph{project-matched} empirical comparison of human-written and LLM-generated Kubernetes configurations. Every LLM receives the same real application's own documentation, and the result is deployed and exercised, not only scanned.
    \item A comparison of five current proprietary and open-weight LLMs. It covers their consistency across repeated samples and their response to an explicit request for production-ready output. This request raises static scores but lowers deployment success.
    \item A root-cause analysis of every configuration that failed at runtime, and a test of whether static security and reliability scores predict deployment success. The test controls for prompt and project.
    \item A replication package containing the docs bundles, prompts, raw model outputs, generated manifests, tool outputs, root-cause labels and analysis scripts.
\end{itemize}

Section~\ref{sec:literature} reviews related work and identifies the research gap. Section~\ref{sec:methodology} describes the dataset, the generation procedure and the benchmark. Section~\ref{sec:results} reports the results for each research question. Section~\ref{sec:discussion} discusses the findings and Section~\ref{sec:threats} the threats to validity. Section~\ref{sec:conclusion} concludes.
